\section{Introduction}

\subsection{Study Rationale}

Pancreatic ductal adenocarcinoma (PDAC) is the third leading cause of cancer
death in the United States and the fourth in the European Union, with total
five-year overall survival below 13 percent \cite{Siegel2025}. The
pancreaticoduodenectomy, or Whipple procedure, is the most technically demanding
of all curative-intent abdominal oncology operations: it requires resection and
reconstruction of four luminal structures (the pancreas, the duodenum, the common
bile duct, and the stomach or pylorus) and dissection around two named vessels
(the superior mesenteric vein and the portal vein), with an artery-first approach
to the superior mesenteric artery. The leading lethal postoperative sequence is
the fistula-sepsis cascade: pancreatic fistula, intra-abdominal infection,
hemorrhage, sepsis, and failure to rescue, graded under the ISGPS classification
of postoperative pancreatic fistula \cite{Bassi2017}. The 2025 Dutch nationwide
cohort of 1000 robotic pancreaticoduodenectomies reported a conversion rate of
10.1 percent, Clavien-Dindo grade III or higher complications of 41.3 percent,
ISGPS grade B or C fistula of 24.4 percent, and 30-day mortality of 3.9 percent
\cite{DutchCohort2025}; these figures anchor the contemporary human baseline. The
present protocol pairs a precise, auditable, on-premises large language model
(LLM) directed robotic Whipple with the RAS(ON) multi-selective inhibitor
daraxonrasib in patients who have completed or are eligible for systemic therapy
such as modified FOLFIRINOX \cite{Conroy2018}.

The clinical argument for the combination is best stated through three
counterfactual scenarios, summarized in Figure 3, in which withholding the
LLM-plus-robot-plus-medicine combination shortens both PFS and OS. In Scenario A,
a borderline-resectable tumor abutting the superior mesenteric vein at 75 degrees
progresses to unresectable disease during a human-only scheduling delay (operating
room access, fatigue, staging lag), and the R0 window closes; the LLM-directed
robot instead achieves a rapid, precise, in-window R0 resection. In Scenario B, an
unrecognized superior-mesenteric or portal-vein injury during manual surgery
triggers the hemorrhage-conversion-fistula-sepsis-failure-to-rescue cascade; the
no-fly vascular gate with a $\leq 3$ ms emergency stop averts the injury and
preserves a complete resection. In Scenario C, manual mistiming of the
daraxonrasib restart (too early or too late) causes anastomotic dehiscence or
micrometastatic outgrowth; the LLM advisory restart at T+7, T+14, or T+21 days,
keyed to fistula status and drug trough, preserves the therapeutic window.

These scenarios invite three framing questions that this protocol is designed to
answer constructively rather than rhetorically. First, given the narrow resection
window and the unforgiving fistula-sepsis cascade, how did we ever rely on humans
alone for this treatment, with no second-opinion oracle isolated from the
operative kinematics. Second, why did we fully trust a prior trial system with so
many opportunities for compounding human error across scheduling, dissection, and
drug timing, when each step admits a documented, deterministic safeguard. Third,
as AI outputs grow more complex and voluminous, why would human peer review alone
continue to suffice for increasingly intricate AI-directed care, rather than a
verification-before-generation discipline with a hash-chained audit trail
\cite{kawchak2026hr9510}. The protocol answers each question by retaining the
human in the loop while making every machine action validated, bounded, and
re-traceable.

\begin{mermaidfig}
\node[mmin]  (a0) at (0,0)      {Borderline PDAC\\tumor abuts SMV at 75 deg};
\node[mmstep](ah) at (0,-2.4)   {Human-only delay\\(OR access, fatigue, lag)};
\node[mmdark](ax) at (0,-4.8)   {Unresectable\\R0 window closes\\shorter PFS/OS};
\node[mmgoal](ac) at (4.6,-2.4) {LLM plus robot\\in-window R0 achieved};
\node[mmin]  (b0) at (9.2,0)    {Intra-op SMV/PV plane};
\node[mmstep](bh) at (9.2,-2.4) {Unrecognized venous injury\\(manual, no zone gate)};
\node[mmdark](bx) at (9.2,-4.8) {Hemorrhage, conversion\\R1/R2, fistula-sepsis};
\node[mmgoal](bc) at (4.6,-4.8) {No-fly gate, $\leq 3$ ms e-stop\\injury averted};
\draw[mmedge] (a0) -- (ah);
\draw[mmedge] (ah) -- (ax);
\draw[mmedge] (a0) -- (ac);
\draw[mmedge] (b0) -- (bh);
\draw[mmedge] (bh) -- (bx);
\draw[mmedge] (b0.west) -- ++(-2.3,0) |- (bc.east);
\end{mermaidfig}
\figcaption{Figure 3. Counterfactual scenarios in which withholding the
combination worsens the patient. Scenario A (left): a human-only scheduling delay
lets a borderline-resectable tumor progress and closes the R0 window. Scenario B
(right): an unrecognized superior-mesenteric or portal-vein injury during manual
surgery triggers the fistula-sepsis cascade; the no-fly gate with a $\leq 3$ ms
emergency stop averts it. A third scenario (drug-restart mistiming) is described
in the text. In each, the LLM-plus-robot-plus-medicine path preserves PFS and OS.}

\subsection{Background}

\subsubsection*{PDAC epidemiology and the Whipple fistula-sepsis cascade}

The epidemiologic burden and operative hazard of PDAC define the unmet need this
protocol addresses. With under-13-percent five-year survival \cite{Siegel2025}
and a most-demanding curative operation whose leading complication is a
self-amplifying fistula-sepsis-hemorrhage cascade \cite{Bassi2017}, the
population has few alternatives and a narrow margin for technical error. The
contemporary robotic baseline \cite{DutchCohort2025} demonstrates that even
expert, high-volume robotic surgery leaves a substantial residual of conversions,
high-grade complications, and clinically relevant fistulae. The combination in
this protocol targets precisely the failure modes that the baseline still
tolerates.

\subsubsection*{Daraxonrasib mechanism and clinical program}

Daraxonrasib (RMC-6236) is an oral RAS(ON) multi-selective inhibitor that binds
the active, GTP-bound (ON) state of multiple RAS variants, providing broad
pan-KRAS coverage including the G12 alleles that predominate in PDAC. The FDA
granted daraxonrasib Breakthrough Therapy designation in June 2025 for previously
treated metastatic PDAC with KRAS G12 mutations \cite{revbreakthrough}, and the
late-phase program comprises RASolute 302 (a Phase 3 study in previously treated
metastatic PDAC) and RASolve 301 (a first-line program) \cite{rasolute302}. The
broad pan-KRAS eligibility of RASolute 302 defines the molecular eligibility used
in this protocol, and the perioperative pause-and-restart logic respects any
RASolve 301 checkpoint-inhibitor combination context. Five prior computational
works by the author independently complemented this drug-development program
throughout 2025 and 2026 and supply the simulation scaffolding adapted here: a
60-second PDAC robotic Whipple and daraxonrasib simulation
\cite{kawchak_2026_20196639}; an AI digital-twin pancreatic-cancer simulation for
in silico FDA-compliance and cost efficiency \cite{kawchak_2025_17239510}; a
quantitative systems pharmacology (QSP) metastatic-PDAC trial simulation with code,
VVUQ, and playbook \cite{kawchak_2025_17001137}; a 100,000-patient in silico
Phase 3 five-arm PDAC trial triplicate \cite{kawchak_2025_16415815}; and
end-to-end PDAC digital-twin clinical-trial proposals \cite{kawchak_2025_15735068}.

\subsubsection*{Current state of AI in clinical trials and surgery}

Artificial intelligence in regulated medicine has, to date, been narrow and
prediction-oriented rather than generative and intervention-directing. As of the
current taxonomy of FDA-authorized medical devices, AI is overwhelmingly deployed
as prediction-only or narrow task-specific software, with reported authorizations
spanning 1,016 in the published taxonomy \cite{Singh2025} and rising to 1,450 or
more by current counts, while zero generative LLM devices have been approved for
patient-facing intervention. This protocol therefore does not propose a generative
device; it proposes a narrow, locked, validated device behavior governed by an
on-premises repository-based LLM whose advisories are bounded, auditable, and
held outside the robot vendor kinematic stack. The author's established record of
LLM evaluation and trust over the period from August 2024 to June 2026 grounds
this design, including a comparative code-generation evaluation of 16 proprietary
versus open-source LLMs against the FDA Adverse Event Reporting System
\cite{kawchak2025codegen16}, a mobile pancreatic-cancer surgical-humanoid program
with priority VVUQ \cite{kawchak2026unitree}, the 21 CFR part 312 Physical AI
unification adaptation \cite{kawchak2026adapt312}, and the
verification-before-generation legislative framing \cite{kawchak2026hr9510}.

\subsubsection*{Physical AI concerns addressed by this protocol}

Eight recurring concerns are raised about Physical AI in patient care, and this
protocol pairs each with a concrete, prespecified mitigation. Figure 4 maps the
eight concerns to their mitigations, and Table 1 states each pairing in full.

\begin{mermaidfig}
\node[mmin]  (c1) at (0,0)      {Physical AI\\limitations};
\node[mmin]  (c2) at (0,-1.9)   {Patient safety\\is paramount};
\node[mmin]  (c3) at (0,-3.8)   {Loss of\\human workers};
\node[mmin]  (c4) at (0,-5.7)   {Single-source\\software};
\node[mmstep](m1) at (5.0,0)     {Phase 0 sim validation;\\narrow locked device};
\node[mmstep](m2) at (5.0,-1.9)  {10 kHz heartbeat,\\$\leq 3$ ms e-stop, force caps};
\node[mmstep](m3) at (5.0,-3.8)  {Human-in-the-loop;\\surgeon plus backup plus ISM};
\node[mmstep](m4) at (5.0,-5.7)  {Cross-framework,\\$\geq 2$ required};
\draw[mmedge] (c1) -- (m1);
\draw[mmedge] (c2) -- (m2);
\draw[mmedge] (c3) -- (m3);
\draw[mmedge] (c4) -- (m4);
\end{mermaidfig}
\figcaption{Figure 4. Representative mapping of Physical AI concerns to protocol
mitigations (first four of eight shown). Each concern, including the proprietary,
non-domestic, complexity, and black-box concerns detailed in Table 1, is paired
with a concrete, prespecified safeguard: validated narrow device behavior over
generative open-endedness, layered hard safety limits, retained human roles, and
multi-framework non-single-source validation.}

\vspace{0.3em}

\noindent\begin{tabularx}{\textwidth}{L{5.0cm} Y}
\toprule
\textbf{Concern} & \textbf{Protocol mitigation} \\
\midrule
Physical AI limitations & Phase 0 simulation validation over $\geq 1000$ procedures and $\geq 2$ independent frameworks; the device behavior is narrow and locked, not generative, and operates only within validated envelopes. \\
Patient safety is paramount & Layered hard limits: a 10 kHz heartbeat broadcast bus, a $\leq 3$ ms cross-arm emergency stop, $\leq 3$ N per-arm and $\leq 18$ N cumulative tip-force caps, and vascular no-fly gating around the SMV and PV. \\
Loss of human workers & The human is retained in the loop: a credentialed surgeon, a backup operator, and an independent safety monitor (ISM) supervise every procedure, with the LLM acting as a bounded advisory oracle rather than a replacement. \\
Single-source software & Validation is cross-framework (for example Isaac, MuJoCo, Gazebo, PyBullet), with $\geq 2$ independent frameworks required, so no single simulator or codebase is a single point of failure. \\
Proprietary-model dependency & The advisory uses an on-premises repository-based LLM with open weights, isolated from any single vendor stack, so the protocol is not captive to one proprietary model provider \cite{kawchak2025codegen16}. \\
Non-domestic open-source LLMs & Deployment is domestic and on-premises, preserving data residency and avoiding reliance on non-domestic open-source model hosting; inference and records remain within the controlled environment. \\
Overly complex AI workflows & A deny-by-default posture with a deliberately simplified gate surface (a small set of Model Context Protocol servers) constrains the action space and reduces workflow complexity to an auditable minimum. \\
LLMs are black boxes & A hash-chained audit trail under 21 CFR part 11 records every advisory, command, and safety event, making each re-traceable to a deterministic seed and a specific commit, so the black box becomes traceable \cite{kawchak2026adapt312}. \\
\bottomrule
\end{tabularx}
\figcaption{Table 1. The eight Physical AI concerns and their protocol
mitigations.}

\subsection{Risk/Benefit Assessment}

\subsubsection{Known Potential Risks}

The known potential risks fall into three categories. Drug risks arise from
daraxonrasib toxicity, including gastrointestinal effects (nausea, diarrhea,
stomatitis), dermatologic and mucocutaneous events, hepatic enzyme elevations,
and dose-limiting toxicities that the 3+3 escalation is specifically designed to
detect and bound; perioperative exposure additionally raises a theoretical risk
to wound and anastomotic healing if restart timing is suboptimal. Device and
robotic risks include device malfunction, an erroneous AI advisory, cyber-physical
failure of the control or telemetry chain, and inadvertent vascular injury; these
are the categories that motivate the significant-risk determination and the
layered hard safety limits. Operative risks are those intrinsic to the Whipple
procedure: postoperative pancreatic fistula (graded by ISGPS \cite{Bassi2017}),
hemorrhage, delayed gastric emptying, intra-abdominal infection and sepsis, and
the need for conversion to open surgery. The contemporary robotic baseline
quantifies the residual magnitude of these operative risks even in expert hands
\cite{DutchCohort2025}.

\subsubsection{Known Potential Benefits}

The known potential benefits are directed at the specific failure modes above. A
precise, in-window R0 resection is the principal oncologic benefit, addressing the
resection-window collapse and the margin-positivity that drive shorter survival.
Layered hard safety limits with a full hash-chained audit trail provide a
benefit not available in manual surgery: every machine action is bounded and
re-traceable, supporting both intraoperative safety and post hoc verification.
Optimized drug-restart timing, advised at T+7, T+14, or T+21 days and keyed to
fistula status and drug trough, aims to restore systemic therapy at the earliest
safe point and thereby suppress micrometastatic outgrowth. For a low-survival
population with few alternatives, even incremental gains in margin status,
complication avoidance, and therapy continuity are clinically meaningful.

\subsubsection{Assessment of Potential Risks and Benefits}

Figure 5 frames the benefit-risk weighing for this population. Against a baseline
of under-13-percent five-year survival and few alternatives, the novel risks
(device malfunction, advisory error, cyber-physical failure) are weighed against
the benefits (precise in-window R0 resection, layered hard safety limits with a
full audit trail, and optimized drug-restart timing). For this low-survival,
limited-alternative population, the benefit-risk balance is favorable: the
potential benefits are likely to outweigh the risks, and the assessment is
consistent with the expanded-access ethic that permits investigational treatment
for serious or immediately life-threatening conditions when the potential benefit
justifies the potential risk (21 CFR \S312.84).

\begin{mermaidfig}
\node[mmdark](base) at (0,0)      {Advanced PDAC baseline\\5-year OS $<$ 13\%\\3rd US / 4th EU death};
\node[mmin]  (r1)  at (-4.6,-2.4) {Device malfunction};
\node[mmin]  (r2)  at (-4.6,-4.4) {AI advisory error};
\node[mmin]  (r3)  at (-4.6,-6.4) {Cyber-physical failure};
\node[mmstep](b1)  at (4.6,-2.4)  {In-window R0 resection};
\node[mmstep](b2)  at (4.6,-4.4)  {Hard limits plus audit trail};
\node[mmstep](b3)  at (4.6,-6.4)  {Optimized restart timing};
\node[mmdec] (w)   at (0,-4.4)     {Benefit-risk for\\low-survival patients};
\node[mmgoal](j)   at (0,-7.4)     {Favorable: benefits likely\\outweigh risks; \S312.84};
\draw[mmedge] (base) -- (w);
\draw[mmedge] (r1.east) -- (w.west);
\draw[mmedge] (r2) -- (w);
\draw[mmedge] (r3.east) -- (w.west);
\draw[mmedge] (b1.west) -- (w.east);
\draw[mmedge] (b2) -- (w);
\draw[mmedge] (b3.west) -- (w.east);
\draw[mmedge] (w) -- (j);
\end{mermaidfig}
\figcaption{Figure 5. Risk-benefit framework for advanced PDAC. Against a baseline
of under-13-percent five-year survival and few alternatives, the novel risks
(device malfunction, advisory error, cyber-physical failure) are weighed against
the benefits (precise in-window R0 resection, layered hard safety limits with a
full audit trail, optimized drug timing). For this low-survival,
limited-alternative population the benefit-risk balance is favorable, consistent
with the expanded-access ethic of 21 CFR \S312.84.}
